% Working revision, 2026-09-20. Original test.tex is intentionally unchanged.
% Evidence and excluded historical results: PAPER_REVISION_NOTES_20260920.md.
% This is an empirical working draft, not a submission-ready boundary claim.
% Use the conference's spconf.sty. References are embedded below.
\documentclass{article}
\usepackage{spconf,amsmath,amssymb,graphicx,booktabs,hyperref}
\newcommand{\U}{\mathbf{U}}
\newcommand{\X}{\mathbf{X}}
\newcommand{\HH}{\mathbf{H}}
\newcommand{\C}{\mathcal{C}}
\newcommand{\R}{\mathbb{R}}
\newcommand{\todo}[1]{\par\noindent\textbf{[TODO]} #1\par}
\newif\ifshowtodos
\showtodostrue % Change to \showtodosfalse only after resolving the TODOs.

\title{RANDOMIZED RECONSTRUCTION OF COMPRESSED MODEL UPDATES:\\
AN EMPIRICAL STUDY OF BIT-WIDTH AND SPARSITY LIMITS}
\name{Anonymous Author(s)}
\address{Anonymous Affiliation}

\begin{document}
\maketitle

\begin{abstract}
How far can communication compression be increased before an update-inversion
attack loses reconstruction quality? We investigate this question using
randomized decoding followed by source separation and feature inversion.
The sender applies deterministic compression; the attacker samples values
inside quantization cells or fills discarded coordinates within observable
top-$k$ bounds. In a fixed 16-image VGG16 experiment, the publicly selected
randomized reconstruction deteriorates from LPIPS 0.5558 at 8 bits to 0.8085
at 4 bits. Randomized decoding does not consistently outperform cell-center
decoding. In a separate, matched 80-image sparsification study, reducing
retention from 50\% to 10\% increases zero-fill LPIPS from 0.4864 to 0.6323.
At 10\% retention, randomized filling yields LPIPS 0.6300, with slightly
lower PSNR and SSIM than zero filling despite a larger search budget.
These measurements characterize degradation under the tested attacks, but
do not yet identify a calibrated image-recovery success/failure boundary
or establish a privacy guarantee.
\end{abstract}

\begin{keywords}
update inversion, randomized decoding, quantization, sparsification,
reconstruction limits
\end{keywords}

\section{Introduction}
Sharing model updates rather than raw data does not by itself prevent
training-data reconstruction. Deep Leakage from Gradients demonstrates
optimization-based reconstruction from gradients~\cite{zhu2019deep}.
The Cocktail Party Attack (CPA) uses independent component analysis (ICA)
to separate the inputs mixed in fully connected layer gradients and
updates~\cite{kariyappa2023cocktail}.
Here we study the behavior of a CPA-based reconstruction pipeline when its
observations are compressed.

Our central question is empirical: \emph{at which bit widths and retention
ratios does a specified attack retain useful reconstruction quality, and
where does it cease to do so?} This is a question about an attack, its
compute budget, and a defined recovery criterion---not a universal claim
that a particular compression setting is secure.

We investigate a simple attacker-side strategy. Instead of using only the
center of each quantization cell or setting every discarded coordinate to
zero, the attacker samples several admissible numerical completions of the
received observation. Each completion is processed by source separation
and feature inversion. Observable scores, rather than private image
metrics, determine which candidates receive further computation.
Randomness is introduced by the attacker, not by the sender's compressor.

The present study provides an 8/7/6/5/4-bit randomized-decoding sweep, a
matched zero-fill retention sweep, and a randomized-filling comparison at
10\% retention. These are complementary experiments, not a completed
factorial study with identical attack backends. We report the observed
quality curves and separate them from the additional evidence needed to
claim a success/failure boundary.

\section{Observation Model}
\label{sec:observation}
\subsection{Multi-step updates with frozen features}
The main experiments use VGG16 with a frozen convolutional feature
extractor and trainable fully connected layers. A client performs $T=10$
full-batch SGD steps on $B=16$ images, using learning rate $\eta=10^{-3}$,
zero momentum, zero weight decay, and disabled dropout. Let
$\HH=[\phi(x_1),\ldots,\phi(x_B)]^\top$ contain the fixed input features
of the first fully connected layer. Writing its per-step gradient as
$\mathbf{A}_t\HH$, with the loss averaging absorbed into $\mathbf{A}_t$,
gives the scaled update
\begin{equation}
  \U=\frac{\mathbf{W}^{(0)}-\mathbf{W}^{(T)}}{\eta}
      =\sum_{t=0}^{T-1}\mathbf{A}_t\HH.
  \label{eq:mixture}
\end{equation}
Thus the frozen-feature setting retains a source-mixture structure across
local steps. Compression perturbs the observed matrix, and neither this
factorization nor successful numerical separation guarantees image recovery.

The attacker observes the compressed update and its decoding metadata,
knows the public model and local training configuration, and attacks the
first fully connected weight matrix. True images, labels, features, and
uncompressed updates are unavailable for candidate construction or ranking.
Private references are used only after selection for evaluation. We do not
assume access to individual updates hidden by an additional secure
aggregation mechanism.

\subsection{Deterministic compression and ambiguous decoding}
For a symmetric signed $b$-bit tensor quantizer, define
\begin{equation}
 Q_b=2^{b-1}-1,\qquad
 s=\frac{\|\U\|_\infty}{Q_b},\qquad
 q_j=\operatorname{round}_{\rm even}(U_j/s).
 \label{eq:quant}
\end{equation}
The scale is stored in FP32, independently for each tensor. A zero tensor
is treated separately. Cell-center decoding returns $s q_j$. The signed
code range is $[-Q_b,Q_b]$, so this convention uses $2^b-1$ numerical levels.
For positive $s$, the ideal interior cell lies between
$s(q_j-\tfrac12)$ and $s(q_j+\tfrac12)$. Actual admissible endpoints also
depend on ties-to-even rounding, finite precision, and the scale-compatible
tensor maximum.

For sparsification, let $r$ denote the \emph{retained fraction}; the removed
fraction is $1-r$. In the matched 80-image study, magnitude top-$k$ is
applied within the attacked first-FC tensor, with
$k=\lfloor rD\rfloor$ for $D$ coordinates. The packet contains a position
bitmap and the retained FP32 values. With retained index set $I$ and
$\tau=\min_{j\in I}|U_j|$, discarded magnitudes satisfy
$|U_j|\leq\tau$. Flat-index order resolves magnitude ties. This is
magnitude-based selection, not random pruning.

\section{Randomized Decoding and Reconstruction}
\label{sec:method}
\subsection{Quantization-cell sampling}
For each received quantized tensor, the attacker samples a completion from
its bounded cells. The implementation uses uniform draws between admissible
FP32 endpoints, then casts to FP32. An extremal coordinate is anchored to
a scale-compatible maximum, and exact re-encoding is checked. Consequently,
the sampled numerical observation is consistent with the received codes;
this does not certify that the final reconstructed images reproduce those
codes after local training. No learned or calibrated gradient prior is used.

We generate eight randomized completions with seeds
$1234+1000003m$, $m=0,\ldots,7$, and retain cell-center decoding as a
separate control. Each first-FC matrix is separated into 16 sources using
FastICA initialization followed by 4,000 gradient-refinement iterations.
The FastICA iteration cap is 200. Source rows are aligned to the
cell-center sources by absolute-cosine assignment, changing only permutation
and sign, not row norms. Absolute source values are then used as targets
for feature inversion (FIA). This backend is a modified CPA-based pipeline,
not an unmodified execution of the original CPA optimizer.

Each randomized candidate receives 5,000 FIA iterations. Its observable
screening score is
\begin{equation}
 J_m=\frac{1}{B}\sum_{i=1}^{B}
 \left[1-\cos\!\left(\phi(\hat x_{m,i}),|\hat h_{m,i}|\right)\right]
       +0.1\,\mathrm{TV}_{\rm norm}(\hat\X_m),
 \label{eq:screen}
\end{equation}
evaluated on a pixel-boxed clone in the model's normalized input space.
The three lowest-scoring candidates continue to a total of 25,000 FIA
iterations, resuming their optimizer and scheduler states. The scheduler
horizon is 25,000 from the first step. Ties use ascending member index.
The rank-1 candidate fixed at screening is the primary randomized result;
we do not replace it with a finalist chosen using private LPIPS.
The center control also receives 25,000 FIA iterations. Total FIA work is
125,000 iterations per bit width including that control, versus 25,000
for the center alone.

\subsection{Bounded random filling for sparsification}
The randomized sparsification experiment uses
\begin{equation}
 \tilde U^{(m)}_j=
 \begin{cases}
   U_j, & j\in I,\\
   \epsilon^{(m)}_j,\quad
     \epsilon^{(m)}_j\sim\operatorname{Unif}(-\tau,\tau), & j\notin I.
 \end{cases}
 \label{eq:fill}
\end{equation}
Retained values remain bit-exact. This bounded completion is an exploratory
sampling model, not an estimate of a validated posterior distribution.
The same five 10\%-retention packets used by the zero-fill control are
attacked with 20 fills per packet. Each fill receives the original CPA
optimizer for 25,000 iterations and FIA for 5,000 iterations.

Unlike the bit sweep, this experiment uses an update-consistency score.
Starting from the public model, the attacker replays the known 10-step
classifier SGD procedure on all 16 reconstructed images, assigning labels
with the public model's predictions. It recompresses that update using the
same top-$k$ codec and minimizes its relative $\ell_2$ discrepancy from
the received decoded packet. The best three candidates continue FIA to
25,000 iterations with preserved optimizer/scheduler states. The same
observable score selects one final, whole 16-image candidate per packet.
This proxy need not select the most accurate private-image reconstruction.

Per packet, the random-fill search uses 500,000 CPA iterations, 160,000 FIA
iterations, and 23 update replays. Zero filling uses 25,000 CPA and 25,000
FIA iterations. Candidate averaging is not used in this experiment, and
these unequal budgets do not isolate the benefit of random filling.

\section{Experimental Results}
\label{sec:results}
\subsection{Evaluation scope}
The quantization sweep uses one fixed 16-image ImageNet batch. The matched
sparsification study uses five fixed 16-image ImageNet batches (80 images),
with the same observations across retention conditions. We report mean
PSNR, SSIM, and LPIPS~\cite{zhang2018lpips}; larger PSNR/SSIM and smaller
LPIPS indicate better scores. Neither image counts nor randomized
candidates are independent repetitions of the complete experiment.

Matching is performed only after the public selection is frozen. The bit
sweep uses pixel-MSE Hungarian matching once per candidate and reuses the
assignment for all image metrics. The 80-image sparsification study uses
signed-source absolute-cosine Hungarian matching after FIA. Because these
studies use different backends and matching protocols, their absolute
metric values must not be interpreted as a controlled comparison between
quantization and sparsification.

\subsection{Bit width: degradation from 8 to 4 bits}
\begin{table}[t]
\centering
\caption{Same 16-image batch, final 25,000-iteration FIA endpoints.
Random denotes the rank-1 candidate selected at 5,000 iterations.
No image-recovery success threshold was registered.}
\label{tab:bits}
\small
\begin{tabular}{clrrr}
\toprule
Bits & Decoding & PSNR $\uparrow$ & SSIM $\uparrow$ & LPIPS $\downarrow$\\
\midrule
8 & Center & 12.0569 & 0.2924 & 0.5594\\
8 & Random & 12.0568 & 0.3070 & 0.5558\\
7 & Center & 11.7955 & 0.2764 & 0.6118\\
7 & Random & 11.8314 & 0.2827 & 0.6332\\
6 & Center & 11.6397 & 0.2596 & 0.6961\\
6 & Random & 11.6675 & 0.2690 & 0.6659\\
5 & Center & 11.5524 & 0.2620 & 0.7826\\
5 & Random & 11.5458 & 0.2558 & 0.7490\\
4 & Center & 11.5197 & 0.2921 & 0.8071\\
4 & Random & 11.5028 & 0.2576 & 0.8085\\
\bottomrule
\end{tabular}
\end{table}

Table~\ref{tab:bits} shows worsening randomized PSNR and LPIPS as bit width
decreases from 8 to 4. Randomized LPIPS increases from 0.5558 to 0.8085,
while PSNR decreases from 12.0568 to 11.5028 dB. SSIM is not monotonic.
Randomization improves LPIPS relative to center decoding at 8, 6, and
5 bits, but worsens it at 7 and 4 bits. At 4 bits the selected randomized
candidate does not improve any of the three reported metrics.

These observations support a quality-degradation account and show that
the tested random search does not consistently compensate for coarser
quantization. They do not establish that 8 bits is successful or that
4 bits is the first failing setting. Completing an optimization run or
passing an observation-provenance check is not an image-recovery criterion.

\subsection{Retention ratio: a matched sparsification sweep}
\begin{table}[t]
\centering
\caption{Same 80 images, original CPA plus FIA. Keep is retained fraction;
remove is sparsity. The four zero-fill rows form the matched retention
sweep. Random filling was evaluated only at 10\% retention in this protocol.}
\label{tab:sparse}
\small
\setlength{\tabcolsep}{3pt}
\begin{tabular}{cc l rrr}
\toprule
Keep & Remove & Fill & PSNR $\uparrow$ & SSIM $\uparrow$ & LPIPS $\downarrow$\\
\midrule
50\% & 50\% & Zero & 22.2878 & 0.3344 & 0.4864\\
30\% & 70\% & Zero & 22.1371 & 0.3141 & 0.5098\\
20\% & 80\% & Zero & 21.8720 & 0.2906 & 0.5536\\
10\% & 90\% & Zero & 21.5305 & 0.2631 & 0.6323\\
10\% & 90\% & Random & 21.5104 & 0.2610 & 0.6300\\
\bottomrule
\end{tabular}
\end{table}

All three zero-fill image metrics worsen as retention decreases
(Table~\ref{tab:sparse}). Importantly, the 50\% packets discard only
already-zero coordinates: these five updates have nonzero fractions
between 0.3822 and 0.4209, and zero-fill decoding reproduces their FP32
values exactly. That condition demonstrates lossless sparse representation
for these observations, not robustness to removing nonzero information.
By contrast, zero-fill relative update errors at 10\% retention range from
0.1141 to 0.1613 across packets.

At 10\% retention, random filling changes mean LPIPS from 0.6323 to
0.6300, while PSNR decreases by approximately 0.0201 dB and SSIM by
0.0021. Without repeated runs or a compute-matched control, this small,
mixed change is not evidence of a reliable reconstruction advantage.
The observed decline across retention levels is not yet a calibrated
transition from image-recovery success to failure.

The actual bitmap-plus-FP32 packets save 46.875\%, 66.875\%, 76.875\%,
and 86.875\% of bytes at 50\%, 30\%, 20\%, and 10\% retention, respectively,
relative to same-shape dense FP32 packets. These figures cover only the
attacked first-FC tensor, not full-model traffic. Random filling adds no
communication because it reuses the received packets.

\section{Discussion and Conclusion}
The completed experiments show a progressive decline in reconstruction
quality with lower bit width and lower retention. Sampling missing
numerical values is a practical way to explore multiple decodings of a
fixed observation, but the tested selectors and budgets do not yield a
consistent advantage over center or zero filling.

The intended boundary statement must therefore remain conditional:
\emph{under a fixed attack budget and a predeclared recovery criterion,
reconstruction meets the criterion at some tested compression levels and
fails to meet it at others}. The present results establish the measured
quality curves, not those final level assignments. A complete randomized
retention sweep and a calibrated per-image criterion remain necessary.
Failure of this attack at a tested setting would still not prove that the
compressed observation prevents other attacks from recovering private data.

\ifshowtodos
\section*{TODOs for the Next Revision (Not Completed Results)}
\todo{T1 --- Define image-recovery success before assigning a boundary.
Specify a held-out calibrated metric or blinded recognition criterion,
its threshold, and the required fraction of recovered images. Report
per-image outcomes, not a threshold applied to mean LPIPS. Distinguish
image success, any-image success per packet, and success per independent
randomized run. Do not infer a boundary below 4 bits merely because the
4-bit program completed.}

\todo{T2 --- Complete a matched randomized retention study. Extend the
same80 protocol to the selected retention levels with the same sampler,
selector, CPA/FIA backend, and budget. The old 80/50/30/10\%-retention
R8 source-averaging runs use a different client arrangement and evaluator;
they cannot fill the missing rows of Table~\ref{tab:sparse}. Prioritize
levels around a transition only after T1. This is a proposed experiment,
not an instruction to launch it automatically.}

\todo{T3 --- Add repeated victim batches and independent complete-search
seeds, with paired uncertainty estimates that respect batch clustering.
Add same-backend FP32 controls where needed and a compute-matched
deterministic-decoding restart control. State both per-candidate and total
CPA/FIA/replay budgets. Do not use image-level historical success rates
as probabilities of independent random trials.}

\todo{T4 --- Unify evaluation for any cross-study comparison. Freeze
matching, preprocessing, pixel range, SSIM settings, and LPIPS version and
backbone. Audit legacy source ordering: existing direct-order and
embedding-matched reevaluations produce different image scores. Keep
private matching strictly evaluation-only.}

\todo{T5 --- Add fixed-index qualitative panels spanning all bit widths
and retention levels, with paired controls and representative failures.
Choose examples without best-of-private-metric selection; include the
complete evaluated set in supplementary material.}

\todo{T6 --- Explain the mechanism with controlled ablations if space
permits: random versus center decoding, candidate count, public screening
quality, and source quality versus image quality. Learned completion
priors, a common update-consistency selector for both codecs, and the
proposed CIFAR-100/FC2 track are extensions, not completed contributions.}

\todo{T7 --- If making a deployment trade-off claim, add measured task
utility and full-model communication accounting. Nominal retention is not
the same as information loss; retain the nonzero-fraction and distortion
audit. No utility or end-to-end privacy guarantee is established here.}

\todo{T8 --- Finish the related-work review with verified primary sources,
resolve all TODOs, supply the conference style, and check the applicable
submission format and page limit. The current file is a working draft.}
\fi

\begin{thebibliography}{9}
\bibitem{zhu2019deep}
L. Zhu, Z. Liu, and S. Han,
``Deep leakage from gradients,''
in \emph{Advances in Neural Information Processing Systems}, vol. 32, 2019.
% https://papers.neurips.cc/paper_files/paper/2019/hash/60a6c4002cc7b29142def8871531281a-Abstract.html

\bibitem{kariyappa2023cocktail}
S. Kariyappa, C. Guo, K. Maeng, W. Xiong, G. E. Suh,
M. K. Qureshi, and H.-H. S. Lee,
``Cocktail Party Attack: Breaking aggregation-based privacy in federated
learning using independent component analysis,''
in \emph{Proc. International Conference on Machine Learning},
vol. 202, pp. 15884--15899, 2023.
% https://proceedings.mlr.press/v202/kariyappa23a.html

\bibitem{zhang2018lpips}
R. Zhang, P. Isola, A. A. Efros, E. Shechtman, and O. Wang,
``The unreasonable effectiveness of deep features as a perceptual metric,''
in \emph{Proc. IEEE Conference on Computer Vision and Pattern Recognition},
2018.
% https://openaccess.thecvf.com/content_cvpr_2018/html/Zhang_The_Unreasonable_Effectiveness_CVPR_2018_paper.html
\end{thebibliography}
\end{document}
